\documentclass[12pt,a4paper]{article}

\usepackage[utf8]{inputenc}
\usepackage[T1]{fontenc}
\usepackage[polish]{babel}
\usepackage{polski}
\usepackage{geometry}
\geometry{margin=2.5cm}
\usepackage{graphicx}
\usepackage{hyperref}
\usepackage{enumitem}
\usepackage{listings}
\usepackage{booktabs}
\usepackage{float}

% ustawienia wygladu kodu
\lstset{
    basicstyle=\ttfamily\small,
    breaklines=true,
    frame=single,
    captionpos=b,
    extendedchars=true,
    keepspaces=true,
    columns=fixed,
    commentstyle=\normalfont\ttfamily,
    showstringspaces=false,
    % zeby polskie znaki w kodzie nie wywalaly bledow
    literate={ą}{{\k{a}}}1 {ć}{{\'c}}1 {ę}{{\k{e}}}1 {ł}{{\l{}}}1 {ń}{{\'n}}1 {ó}{{\'o}}1 {ś}{{\'s}}1 {ź}{{\'z}}1 {ż}{{\.z}}1 {Ą}{{\k{A}}}1 {Ć}{{\'C}}1 {Ę}{{\k{E}}}1 {Ł}{{\L{}}}1 {Ń}{{\'N}}1 {Ó}{{\'O}}1 {Ś}{{\'S}}1 {Ź}{{\'Z}}1 {Ż}{{\.Z}}1
}

% informacje o autorach
\title{Dokumentacja Funkcjonalna Projektu C\\ \large System wyznaczania współrzędnych węzłów grafów planarnych}
\author{Autorzy: Mateusz Bombola, Kajetan Karwacki}
\date{\today}

\begin{document}

\maketitle
\newpage
\tableofcontents
\newpage

\section{Cel projektu}
Celem projektu jest stworzenie aplikacji w języku C, która na podstawie topologii grafu planarnego wyznacza optymalne współrzędne dwuwymiarowe $(x, y)$ dla jego wierzchołków. Program służy jako silnik obliczeniowy przetwarzający dane o połączeniach w celu wygenerowania układu umożliwiającego czytelną wizualizację grafu.

\section{Opis algorytmów rozmieszczania grafu}
W programie zaimplementowano dwie niezależne metody wyznaczania współrzędnych:

\begin{itemize}
    \item \textbf{Algorytm SMACOF (\texttt{smacof}):} Służy do rozmieszczania elementów w przestrzeni na podstawie ich wzajemnych podobieństw lub odległości. Działa poprzez iteracyjne poprawianie położenia punktów tak, aby odległości między nimi jak najlepiej odpowiadały zadanym wartościom. W efekcie powstaje układ, w którym elementy bardziej podobne znajdują się bliżej siebie, a mniej podobne dalej.
    
    \item \textbf{Algorytm Fruchtermana-Reingolda (\texttt{fr}):} Automatycznie rozmieszcza elementy grafu w przestrzeni tak, aby był on czytelny i przejrzysty. Traktuje węzły jak obiekty, które wzajemnie się odpychają, a połączone elementy przyciągają się jak sprężyny. Dzięki wielokrotnemu przeliczaniu tych zależności program stopniowo ustala układ, w którym struktura grafu jest łatwa do zrozumienia.
\end{itemize}

\section{Interfejs linii poleceń}
Program jest sterowany poprzez argumenty przekazywane w wierszu poleceń. Składnia wywołania:
\begin{center}
    \begin{verbatim}
    ./graph_layout -i <input> -o <output> -a <smacof|fr> [-n iter] [-f txt|bin]
    \end{verbatim}
\end{center}

\subsection{Lista dostępnych flag}
\begin{table}[H]
\centering
\begin{tabular}{@{}lll@{}}
\toprule
\textbf{Flaga} & \textbf{Argument} & \textbf{Opis} \\ \midrule
\texttt{-i} & \textit{ścieżka} & Ścieżka do wejściowego pliku tekstowego (wymagane). \\
\texttt{-o} & \textit{ścieżka} & Ścieżka do pliku wynikowego (wymagane). \\
\texttt{-a} & \textit{algorytm} & Wybór algorytmu: \texttt{smacof} lub \texttt{fr}. \\
\texttt{-n} & \textit{liczba} & Liczba iteracji (domyślnie 1000). \\
\texttt{-f} & \textit{format} & Format wyjściowy: \texttt{txt} lub \texttt{bin}. \\
\texttt{-h} & --- & Wyświetlenie pomocy technicznej. \\
\bottomrule
\end{tabular}
\caption{Parametry sterujące programem}
\end{table}

\section{Formaty danych}
\subsection{Format wejściowy}
Każda linia reprezentuje krawędź: \texttt{<nazwa> <od\_id> <do\_id> <waga>}.
\\
\textbf{Przykład:}
\begin{lstlisting}
AB   1  2  1
BC   2  3  1
CD   3  4  1
DB   4  2  1.407
\end{lstlisting}

\subsection{Format wyjściowy}
Wynikiem działania programu ma być plik tekstowy lub binarny (w zależności od wyboru użytkownika) zawierający listę współrzędnych wierzchołków: \texttt{<wierzchołek> <współrzędna\_x> <współrzędna\_y>}.
\\
\textbf{Przykład:}
\begin{lstlisting}
1 0.0 0.0
2 1.0 0.0
3 1.0 1.0
4 0.0 1.0
\end{lstlisting}

\section{Kody błędów}
\begin{itemize}
    \item \textbf{Kod 1:} Błąd otwarcia pliku wejściowego.
    \item \textbf{Kod 2:} Nieprawidłowy format danych wejściowych.
    \item \textbf{Kod 3:} Nieznany algorytm (oczekiwano \texttt{smacof} lub \texttt{fr}).
    \item \textbf{Kod 4:} Błąd alokacji pamięci.
\end{itemize}

\section{Przykłady użycia}
% gotowe komendy do skopiowania

1. Wyznaczenie współrzędnych metodą Fruchtermana-Reingolda i zapis do tekstu:
\begin{lstlisting}[language=bash]
./graph_layout -i dane.txt -o wynik.txt -a fr -n 1000 -f txt
\end{lstlisting}

2. Wyznaczenie współrzędnych metodą SMACOF i zapis binarny (500 iteracji):
\begin{lstlisting}[language=bash]
./graph_layout -i siec.txt -o mapa.bin -a smacof -n 500 -f bin
\end{lstlisting}

\end{document}
